\section{Executive Accountability: responsabilidad compartida, pero no diluida}

Los incidentes de seguridad graves suelen dejar al descubierto la estructura de la organización: quién recibe los hechos, quién puede exigir más investigación, quién aprueba los pagos, quién decide no notificar, quién informa al consejo de administración. Cargar toda la responsabilidad sobre el CSO fomenta una cultura defensiva. Repartirla por igual entre una decena de comités, en cambio, hace que nadie sea realmente accountable. Una gobernanza madura necesita que coexistan la shared responsibility y un named decision owner.

\subsection{Los decision rights del CSO, el GC, el CEO y el board}

El CSO debe asegurarse de que el riesgo se escale con precisión. El GC debe asegurarse de que un equipo calificado atienda las cuestiones legales. El CEO debe responder por los recursos y por las decisiones de negocio importantes. El board o el disclosure committee debe aportar supervisión y challenge. El CSO no puede sustituir con un «legal approved» su propia obligación de veracidad técnica. Tampoco el GC puede juzgar el impacto por su cuenta cuando falta evidence. Y menos aún debe el CEO oír hablar por primera vez de la arquitectura de seguridad justo antes de la divulgación pública.

\lecturefigure{12-executive-accountability.png}{La responsabilidad en ciberseguridad la comparten el CSO, el GC, el CEO y el board, pero cada decisión necesita un owner explícito.}{Subtítulos locales 25:20--35:20; conceptos de gobernanza redibujados.}

\begin{knowledgebox}{Lectura de la figura: se comparte la información, no una responsabilidad difusa}
El CSO aporta el technical scope, la security recommendation y los escalation facts. El GC aporta los legal duties, la regulator strategy y el privilege advice. El CEO toma las decisiones de enterprise risk y de disclosure. El board o el committee supervisa, cuestiona y exige una revisión posterior. Los límites de la figura no son silos: cada rol puede hacer challenge a los demás workstreams, pero no puede trasladar a otro su propio deliverable. El registro final debe poder responder quién es responsable de los hechos, quién del análisis y quién aprobó las acciones.
\end{knowledgebox}

\teachervoice{El ponente recomienda que los líderes de seguridad construyan la relación con el CEO y el GC antes de una crisis, en lugar de pedirles que entiendan un incident complejo en la primera llamada importante. Nota de clase: la executive readiness es una interface que se entrena de antemano y requiere brief periódicos, tabletop y un umbral explícito de «cuándo despertarme».}

\begin{warningbox}{«Ya se lo había dicho» no significa que la gobernanza esté completa}
Un solo correo o un aviso verbal no demuestra que el destinatario entienda el riesgo, tenga el decision right o haya completado las acciones siguientes. Un escalamiento de alto riesgo debe incluir severity, evidence, unknowns, recommended actions, deadline, un ask explícito y una confirmación de recepción. Si la decisión se aplaza o se rechaza, deben registrarse el owner, el motivo y la siguiente revisión.
\end{warningbox}

\begin{importantbox}{La forma correcta de proteger a los líderes de seguridad}
Proteger no significa redactar memos para cubrirse. La organización debe proporcionar una job authority clara, un escalamiento por escrito, asesoría legal independiente, acceso al consejo de administración, D\&O/indemnification y el employment support adecuado. Por su parte, el líder debe mantenerse honesto respecto de los hechos, no excederse en sus atribuciones emitiendo conclusiones legales, negarse a eliminar registros desfavorables y, ante desacuerdos importantes, recurrir a los canales formales de escalamiento.
\end{importantbox}

\subsection{Resumen de la sección}

El objetivo de la accountability es que decidan las personas que tienen la capacidad, la información y la autoridad para hacerlo. El CSO no debe convertirse en el único responsable de todos los riesgos de seguridad, pero tampoco puede renunciar, con el pretexto de una «decisión colectiva», a escalar con precisión ni a su juicio profesional.
